\documentclass[11pt]{article}
\usepackage[T1]{fontenc}
\usepackage{lmodern}
\usepackage[margin=1in]{geometry}
\usepackage{booktabs}
\usepackage{array}
\usepackage{xcolor}
\usepackage[colorlinks=true,allcolors=blue!60!black]{hyperref}
\usepackage{parskip}

\title{Can AI Startups Replace the Jobs Big Tech Cuts?\\[0.3em]\large I Ran 2,000 Simulations to Find Out}
\author{Gerard Spivey\\\href{mailto:gs@gspivey.com}{gs@gspivey.com}}
\date{October 2026}

\begin{document}
\maketitle

\textit{The short version of my working paper, ``Fewer Giants, More Firms?''}

\hrule
\vspace{1em}

\section*{The Question I Couldn't Shake}

On October 26, Microsoft offered a voluntary exit package to about \textbf{60,000 employees}, roughly a quarter of its workforce, while guiding to around \textbf{\$175 billion} in 2026 capital spending as it builds out AI data centers.

AI already writes most of my implementation code, so the ``AI is coming for jobs'' headlines aren't abstract for me. But they all assume a fixed set of companies, each needing fewer people.

That misses half the story. If AI makes big companies more productive, it should also make \textit{new} companies cheaper to start. So I asked: \textbf{when big companies shed jobs because of AI, will new AI-powered companies hire enough people to make up the difference?} Economists call that \textbf{offset}.

\section*{Why Cloud Computing Is the Warning Sign}

The closest precedent is \textbf{cloud computing}. After AWS, you no longer had to buy servers to launch a software company. Starting got dramatically cheaper.

You'd expect a startup boom. Instead, between 2006 and 2020, \textbf{the US startup rate kept falling and employment kept concentrating in big firms.} Cheaper to start didn't mean more jobs.

To find out what \textit{does} decide it, I built a model of an economy where companies are born, grow, and die, then hit it with an AI shock that both cuts jobs at incumbents and lowers the cost of starting a firm. Nobody knows how strong those effects will be, so I ran it \textbf{2,000 times} with different plausible assumptions. That's a \textbf{Monte Carlo simulation}: instead of betting on one forecast, you look at the whole spread of outcomes.

\section*{The 81.7\% Headline (and Why It's Not the Story)}

In \textbf{81.7\%} of those simulated futures, AI raised total employment. Good news? Partly. When I broke down what drove the results, one factor dominated:

\begin{center}
\begin{tabular}{p{7cm}r}
\toprule
\textbf{What drives the outcome} & \textbf{Share of variation} \\
\midrule
Whether AI \textbf{helps} workers or \textbf{replaces} them & \textbf{71\%} \\
How much AI lowers the cost of starting a company & \textbf{3\%} \\
\bottomrule
\end{tabular}
\end{center}

Cheaper startups, the cloud-era mechanism, barely move the needle. Most of that 81.7\% comes from worlds where AI simply \textbf{isn't replacing many people to begin with}, not from startups riding to the rescue.

\section*{The Tipping Point}

Think of AI as sitting on a dial between \textbf{power tool} and \textbf{replacement}. As a power tool, it makes each worker more productive, adopting companies grow, and growing companies hire. As a replacement, companies make the same product with fewer people.

Somewhere on that dial is a \textbf{tipping point}, which I call the \textbf{displacement threshold}. Which side we land on changes everything:

\begin{center}
\begin{tabular}{p{8cm}r}
\toprule
\textbf{Where AI lands} & \textbf{Chance AI adds jobs} \\
\midrule
\textbf{Power-tool side} of the tipping point & \textbf{99.8\%} \\
\textbf{Replacement side} of the tipping point & \textbf{74.9\%} \\
Far past it, and AI-native startups aren't much more productive & \textbf{27.8\%} \\
\bottomrule
\end{tabular}
\end{center}

Past the tipping point, the deciding question becomes: \textbf{are AI-native startups dramatically more productive?} That matters about six times more than how cheap they are to start.

\section*{Where We Are Today}

The early data is cautiously encouraging. The Census Bureau's October 2026 survey found \textbf{only 2\% of businesses using AI cut staff because of it}. That's the power-tool side, for now.

The worrying part: \textbf{new AI companies barely hire.} By press reports, Cursor brings in about \textbf{\$10 million in revenue per employee}. Cloud-era Snowflake and Databricks were at \$0.17--0.32 million. Business applications are at a record, up \textbf{15.3\%} this year, but applications from people planning to \textbf{pay wages} are \textit{down} 15.3\%. That decline predates AI and shows up in every sector, but the direction is clear: \textbf{more firms, not more employers.}

\section*{US vs. Europe}

I also ran a European version of the model:

\begin{center}
\begin{tabular}{p{6.5cm}rr}
\toprule
& \textbf{United States} & \textbf{European Union} \\
\midrule
New employer firms (share/year) & 9.1\% & 4.8\% \\
Private AI investment (2025) & \$285.9B & \$20.9B \\
Paperwork to start (OECD, 0--6) & 2.03 & 1.11--2.74 \\
Firing protection (OECD, 0--6) & 0.77 & 1.94--2.88 \\
AI-specific regulation & None federally & AI Act (Dec 2027) \\
\textbf{Simulated chance AI adds jobs} & \textbf{81.7\%} & \textbf{69.7\%} \\
\bottomrule
\end{tabular}
\end{center}

The surprise: \textbf{Europe's problem isn't startup paperwork.} The US actually scores worse than France, Italy, Spain, and Sweden there. Most of the gap is \textbf{thin scale-up money}. Europe has AI startups, but far less capital to grow them into real employers.

\section*{What I'm Watching}

These leading indicators will tell us which side of the tipping point we're on:

\begin{center}
\begin{tabular}{p{4.5cm}p{4cm}p{3.5cm}}
\toprule
\textbf{Signal} & \textbf{Where to find it} & \textbf{Bad sign} \\
\midrule
AI-using businesses cutting staff & Census BTOS (US), Eurostat (EU) & Climbs from 2\% \\
Applications planning to pay wages & Census BFS & Keeps falling \\
Headcount at AI-native startups & Census/OECD firm data & Revenue grows, headcount flat \\
\bottomrule
\end{tabular}
\end{center}

\textbf{The first row is the big one.} If that 2\% starts climbing, the 99.8\% world is over.

\section*{My Takeaway}

Honest caveats first: this is a model, not a forecast. It counts jobs, not whether they pay well or whether a laid-off worker can land one, and cloud computing is an imperfect stand-in for AI.

With that said: \textbf{cheaper startups won't save us on their own.} What decides AI's effect on jobs is whether it stays a power tool or becomes a replacement, and right now the data says power tool. I'll check the Census survey each time it updates, and if that 2\% moves, I'll write the follow-up.

\vspace{1em}
\hrule
\vspace{0.5em}
\textit{Want the full paper and the model behind it?} \href{https://www.gspivey.com/contact}{Reach out.}

\end{document}
